\begin{lemma}[Stable change, Cases 1 and 2, with a transported split tail]
Fix a finite loopless multigraph with its individual edge type, a palette
of \(k\) colors, a distinguished edge \(e_0\), an input coloring, an
ordered tree-sequence \(T\), a generating coloring \(\varphi_n\), a
positive-rung ETT \(\mathcal E\) on \(T\), a hierarchy base, and a
\noderef{HierarchyCounterexample} \(C\).  Assume the properness,
exactly-one-uncoloredness, complete lower-rung six-pack, fewer-divider
induction, rung-local lexicographic minimality, and \(2n+11\) estimate
recorded in a
\noderef{HierarchyMinimumCounterexampleContext}.

Write \(H=C.\mathit{hierarchy}\), let \(q\) be its last stored index, and
put
\[
 A=\noderef{HierarchyStarredStage}(H,q),\quad
 B=\noderef{HierarchyFullTailPrefix}
       (C,\noderef{HierarchyFullTailLength}(C)-1).
\]
Let \(\alpha\ne\beta\) be missing on \(B\).  Let \(z\) certify a possibly
singleton \noderef{TwoColorPathComponent} \(Q\) under \(\varphi_n\), put
\[
 \sigma_n=\noderef{TwoColorComponentSwap}
      (G,\varphi_n,\alpha,\beta,z),
\]
and assume that \(Q\) is disjoint from the full component through the
first displayed vertex of \(T\) missing \(\alpha\).  Assume the explicit
Cases 1/2 alternative: either \(q>0\) and \(\alpha\) is missing on
\(T_{n,q}\), or \(q=0\) and either \(\alpha\) is missing on \(T_n\), or
\(\alpha\) is missing on \(A\) and both \(\alpha,\beta\) lie outside
\(D_n=\noderef{TashkinovStoredColors}(\mathcal E,n)\).

Let \(X\) be an exact split prefix satisfying
\[
             A\prec X\preccurlyeq B
\]
in the sense of \noderef{TreeSequenceInitialSegment}.  Assume the
following transport data.
\begin{enumerate}
\item Every individual displayed edge of \(T\) at a position \(r\) with
\(|A|\le r<|X|\) keeps its value under \(\sigma_n\).
\item There is a verified \(\sigma_n\)-augmentation route through \(X\):
if the last extension kind is PE, it is a
\noderef{TashkinovAugmentationSequence} \(T_n\longrightarrow X\); if the
kind is SE or RE, there are the exact edge \(f_n\), exact seed
\(T_n+f_n\), singleton equality \(F_n=\{f_n\}\), and a
\(\sigma_n\)-augmentation sequence from that seed to \(X\).
\item There is an original-color augmentation suffix
\(X\longrightarrow T\), and both \(\alpha\) and \(\beta\) are missing on
\(X\) under \(\sigma_n\).
\item For every vertex \(v\in V(B)\),
\(\overline\sigma_n(v)=\overline\varphi_n(v)\).
\item For every \(\eta\in D_{n,q}=\noderef{HierarchyStoredColorsAt}(H,q)\),
the totalized first missing position on \(T\) is the same under the two
colorings.
\item If either affected color belongs to the copied last Gamma pair
\(C.\mathit{good}.\Gamma^q_\eta\), then the old first missing position of
\(\eta\) is at most \(|X|\).
\end{enumerate}

Then \(\sigma_n\) is \noderef{HierarchyWeaklyStable} at \(A\) and is a
\noderef{ModColoring} of \(\mathcal E\).  Moreover there exist an actual
\noderef{CorrespondingExtendedTashkinovTree} certificate on the same full
ordered tree \(T\), maximum property for its replacement ETT, a hierarchy
base for that ETT, a full hierarchy \(K\), and a
\noderef{GoodHierarchy} \(g_K\), such that
\[
                 \noderef{SameGoodHierarchy}(C.\mathit{good},g_K).
\]
Thus the conclusion constructs, rather than assumes, the full-tree ETT,
maximum property, base, hierarchy, goodness, and same-good data; it makes
no assertion that all affected tail-edge values are fixed.
\end{lemma}
\begin{proof}
Let \(T_n\) be the terminal ladder tree and retain the notation
\(H,q,A,B,Q\) from the statement.  First derive the fact (6.6) which is
not a field of the minimum context.  Apply
\noderef{HierarchyCounterexamplePenultimateElementary}, in order, to
\(G,k,e_0\), the input coloring, \(T\), \(\varphi_n\), the present ETT,
base, and counterexample, followed by exactly
\[
 \mathit{ctx.proper},\quad \mathit{ctx.oneEdgeUncolored},\quad
 \mathit{ctx.lowerRung},\quad \mathit{ctx.fewerDividers},\quad
 \mathit{ctx.lexicographicMinimality},\quad
 \mathit{ctx.largeMissingSet}.
\]
The six displayed arguments are respectively the properness,
exactly-one-uncoloredness, complete lower-rung six-pack, fewer-divider
induction, rung-local lexicographic minimality, and \(2n+11\) projections
of the stated \noderef{HierarchyMinimumCounterexampleContext}.  By the
definitions of \noderef{HierarchyFullTailLength} and
\noderef{HierarchyFullTailPrefix}, the exact conclusion of the cited
lemma is
\[
 h_B:\quad
 V(B)\text{ is elementary under }\varphi_n.                      \tag{0}
\]
Thus \(h_B\) is derived from those six projections; it is not itself a
context projection.

Now apply \noderef{StableChangeCasesOneTwo} to the same
\(G,k,e_0\), input coloring, \(T\), \(\varphi_n\), ETT, base, and
counterexample; pass, in its displayed order, the six context projections
above with the newly derived fact \(h_B\) between
\(\mathit{ctx.lexicographicMinimality}\) and
\(\mathit{ctx.largeMissingSet}\); and then pass \(\alpha\ne\beta\), the
two supplied colors missing on \(B\), the path certificate for \(Q\), the
defining component-swap equality for \(\sigma_n\), the Cases 1/2
alternative, and the component disjointness.  We obtain
\[
 \tag{1}
 \noderef{HierarchyWeaklyStable}(H,q,\sigma_n),\qquad
 \noderef{ModColoring}(\mathcal E,\sigma_n),
\]
and, when \(q>0\), an actual maximum ETT on \(A\) with a good hierarchy
through the stage preceding \(q\), the old core and closure, all retained
\(D\)-sets and Gamma pairs, and the required last minus-closed assertion.
These are conclusions of the cited lemma, not additional premises here.

We next construct the exact full-tree route.  The original suffix is a
\(\varphi_n\)-augmentation sequence from \(X\) to \(T\), while both
affected colors are missing on \(X\) under the displayed component swap.
Therefore
\noderef{TashkinovAugmentationSequenceComponentSwap}, applied with the
same root \(z\), gives the same exact ordered sequence and individual
new-edge identities as a \(\sigma_n\)-augmentation sequence
\[
                         X\longrightarrow T.                 \tag{2}
\]
This application is precisely why no assertion about the colors of the
suffix edges is needed.

If the final kind is PE, concatenate the supplied
\(\sigma_n\)-route \(T_n\longrightarrow X\) with (2).  Transitivity of the
reflexive-transitive relation in
\noderef{TashkinovAugmentationSequence} gives
\[
                         T_n\longrightarrow T.               \tag{3P}
\]
If the final kind is SE or RE, retain the supplied individual identity
\(f_n\), the literal equality \(F_n=\{f_n\}\), and the supplied
\noderef{TreeSequenceAdjoinEdge} certificate for the exact seed.  Concatenate
its route \(T_n+f_n\longrightarrow X\) with (2), obtaining
\[
                         T_n+f_n\longrightarrow T.           \tag{3N}
\]
Thus (3P) and (3N) have exactly the two shapes required by the universal
clause of \noderef{ModColoring}; no seed is reconstructed from endpoint
data, and parallel edges remain distinguished.

The `abovePenultimate' and `strictAbovePenultimate' fields in the
definition \noderef{ExtendedTashkinovTree} say, for this old ETT, exactly
that \(T_n\preccurlyeq T\) and \(|T_n|<|T|\).
Instantiate the second conclusion in (1) at the actual tree \(T\) and
the relevant route (3P) or (3N).  It returns a nonempty
\noderef{CorrespondingExtendedTashkinovTree}; choose a certificate
\(\mathcal C_T\).  Applying \noderef{MaximumPropertyTransport} to (1),
\(\mathcal C_T\), and the maximum-property field of \(C\) proves that
\(\mathcal C_T.\mathit{replacementETT}\) has maximum property.          \(\tag{4}\)

It remains to build the base, hierarchy, and goodness on this actual
replacement ETT.  First suppose the final kind is PE.  Unfolding the
weak stability in (1) gives the required
\noderef{CombinedStableColoring} on \(T_n\) and the PE closure, and (3P)
is the required full augmentation route.  Apply
\noderef{HierarchySplitter} with the lower-rung field of the context.
It yields a corresponding certificate (which we now use as
\(\mathcal C_T\)), its transported maximum property, a replacement base
with the old core and closure, and a hierarchy \(K\) with exactly the old
ordered stages and divider positions.  Thus the PE construction does not
identify two independently chosen existential certificates.

If the final kind is SE or RE, use \(\mathcal C_T\) from (4) directly.
The fields of \noderef{CorrespondingExtendedTashkinovTree} identify the
rung, every positive ladder tree, every raw stored set, the exact
connecting-edge set, and the extension kind.  Define the replacement
base core to be the replacement terminal ladder tree.  The old
\noderef{HierarchyBase}.nonparallel clause identifies the old core with
\(T_n\), so this is the same exact ordered tree after the ladder-tree
identification.  Copy the numerical support position and unused closure;
the terminal missing-color clause is reflexive under \(\sigma_n\), the PE
implication is vacuous, and the non-PE clause is the defining core
equality.  Copy the divider count, stages, and positions of \(H\).
All fields of \noderef{TashkinovHierarchy} are coloring-free exact-prefix,
strict-length, and exact-entry assertions, so they construct the required
hierarchy \(K\) on the replacement ETT.

We now compare the stored sets.  Correspondence identifies the raw
stored arrays.  By \noderef{TashkinovStoredColors}, the computed \(D_n\)
is that common bounded union with the terminal-ladder missing colors
removed.  Weak stability in (1) identifies the missing set at each vertex
of \(T_n\), hence identifies the two computed copies of \(D_n\).
Every starred stage is an exact prefix of \(A\), and
\(A\preccurlyeq X\preccurlyeq B\).  For a vertex of a starred stage,
unfold \noderef{TreeSequenceVertexSet}, carry its displayed index through
these exact \noderef{TreeSequenceInitialSegment}s, and apply the stated
vertexwise equality on \(B\).  Thus the two colorings have equal
tree-missing sets on every starred stage.  Unfolding
\noderef{HierarchyStoredColorsAt} gives, at every index \(j\),
\[
                         D^K_{n,j}=D^H_{n,j}.                 \tag{5}
\]
The identical index argument on \(T_n\) for \(j=0\), and on the retained
positive hierarchy stages for \(j>0\), gives equality of the two
\noderef{HierarchyStageMissingColors} sets as well.             \(\tag{6}\)

Define the new Gamma pair by copying
\(C.\mathit{good}.\Gamma^j_\eta\), with the canonical index casts from
the equal divider counts.  Equations (5)--(6) immediately transport the
two-color cardinality and pairwise-disjointness fields of
\noderef{GoodHierarchy}.  Whenever \(q>0\), independently of the terminal
extension kind, the prefix good hierarchy supplied by
\noderef{StableChangeCasesOneTwo} transports the parallel-base field
when applicable and all unique-refresh and minus-closed fields before
the last divider.  Its separate last minus-closed conclusion supplies the
last such field.  The last unique-refresh field itself quantifies only
vertices of stages contained in \(A\preccurlyeq B\), so (5), (6), and the
vertexwise equality on \(B\) transport the old unique witness and its
uniqueness.  When \(q=0\), there are no refresh or minus-closed fields,
and in the SE/RE case the parallel-base implication is vacuous.

For completeness, consider the PE parallel-base field when \(q=0\).
Here \(A\) is the ordered-union core.  Since \(A\prec X\), the supplied
route through \(X\), whose steps increase length by one, has an exact
intermediate stage \(A\); split it there and unpack the first
\(\sigma_n\)-\noderef{TashkinovAugmentation} above \(A\).  It supplies
the exact boundary edge \(g=T.\mathit{edge}(|A|)\), its
\(\sigma_n\)-color \(c\), and a \(\sigma_n\)-missing witness for \(c\)
on \(A\).  The first transport assumption fixes this individual edge
because \(|A|<|X|\), so its \(\varphi_n\)-color is also \(c\).
Vertexwise equality on \(B\), restricted along \(A\preccurlyeq B\),
turns the same witness into a \(\varphi_n\)-missing witness.  The attachment
field of the tree sequence makes \(g\) a literal
\noderef{BoundaryEdges} member of \(V(A)\), so it is a
\noderef{BoundaryColorEdges} member of color \(c\) under \(\sigma_n\).
If \(c\) lay in the stage-zero Gamma union, old
\noderef{GoodHierarchy}.avoidance and the definition of
\noderef{HierarchySegmentColors} would forbid exactly this first edge:
membership of \(\eta\) in
\(\noderef{HierarchyStoredColorsAt}(H,0)\) excludes \(\eta\) from the
colors missing on \(A\), so its first missing position is after
\(|A|\), while the old value of this edge is \(c\).  Thus \(c\) is unreserved and
the core is not \noderef{TreeColorSetMinusClosed} under \(\sigma_n\).
The two stable clauses of the combined stability identify missing sets
on \(T_n\) and on the PE closure; together with (5) and the copied Gamma
pairs, they transport the old cardinality-four clause.  This proves both
parts of the parallel-base field without asserting equality of all
boundary-edge colors.

Only avoidance at the last stored index remains.  Fix
\(\eta\in D^H_{n,q}\) and
\(c\in C.\mathit{good}.\Gamma^q_\eta\).  Old avoidance says that \(c\)
is stage-missing and that no exact edge position \(r\) satisfies
\[
 |A|\le r<
 \min\{\noderef{HierarchyFirstMissingPosition}
          (\varphi_n,T,\eta),|T|\},
 \qquad \varphi_n(T.\mathit{edge}(r))=c.                    \tag{7}
\]
The stage-missing part transports by (6), and the endpoint in (7) is
unchanged by the stored-first-position hypothesis.

If \(c\notin\{\alpha,\beta\}\), the definition of
\noderef{TwoColorComponentSwap} fixes every edge whose old or new value
is \(c\).  A new segment witness under \(\sigma_n\) would therefore be
the identical old witness forbidden by (7).  If
\(c\in\{\alpha,\beta\}\), the Gamma-cutoff premise says that the common
first missing position of \(\eta\) is at most \(|X|\).  Every candidate
\(r\) for the new segment then satisfies
\(|A|\le r<|X|\), where the explicit fixed-prefix premise gives
\(\sigma_n(T.\mathit{edge}(r))=\varphi_n(T.\mathit{edge}(r))\).
Again a new witness is an old witness forbidden by (7).  This checks the
last avoidance clause pointwise.  It neither equates the full segment
color sets nor fixes any affected edge after \(X\).

All fields of a \noderef{GoodHierarchy} \(g_K\) on \(K\) have now been
constructed.  The replacement base has the same core and PE closure,
the hierarchy has the same divider count, ordered stages, and positions,
and (5) together with the definition of \(g_K\) gives equal stored sets
and Gamma pairs at every index.  These are exactly the clauses of
\noderef{SameGoodHierarchy}(C.\mathit{good},g_K).  Combining this with
(1) and (4) proves the stated conclusion.
\end{proof}
